\documentclass[../../main.tex]{subfiles}
\begin{document}

{\bf [解]}
\paragraph{(1)}

簡単さのため$q_k=f_n(kh)$ とおくと，題意の方程式は
\begin{align}
  \begin{cases}
    q_0 = c \\
    \dfrac{q_{k+1}-q_k}{h} = (1-q_k)q_{k+1}
  \end{cases} \label{eq:1}
\end{align}
である．

まず任意の$k$について$q_{k}\neq 0$であることを背理法で示す．
$k \ge 0$に対して$q_{k}=0$と仮定する．\cref{eq:1}の第二式で$q_{k+1}=0$を代入すると$q_{k}=0$となり，帰納的に$q_0=0$となるが，これは\cref{eq:1}で$q_0=c \neq 0$に反して矛盾である．よって背理法により
\begin{align}
  q_{k} \neq 0 \label{eq:2}
\end{align}
である．

そこで\cref{eq:1}の第二式の両辺を $q_k q_{k+1} (\neq 0)$ で割って，$p_k=\dfrac{1}{q_k}$より
\begin{align}
  \dfrac{p_k - p_{k+1}}{h} = p_k - 1 \quad \therefore p_{k+1} - 1 = (1-h)(p_k - 1)
\end{align}
だから，くり返し用いて $p_0 = \dfrac{1}{q_0}=\dfrac{1}{c}$ より，
\begin{align}
  p_k &= \dfrac{1-c}{c} (1-h)^k + 1 \label{eq:3}
\end{align}
である．

\paragraph{(2)}

\cref{eq:3}で$k=n$とすると$kh=nh=n\dfrac{a}{n}=a$だから
\begin{align}
  p_n &= \dfrac{1}{f_n(a)}=\dfrac{1-c}{c} (1-h)^n + 1 \\
  \therefore
  f_n(a) = \dfrac{1}{\dfrac{1-c}{c} (1-h)^n + 1}
\end{align}
である．$h=\dfrac{a}{n}$を代入して
\begin{align}
  f_n(a)
  &= \dfrac{1}{\dfrac{1-c}{c} \left(1-\dfrac{a}{n}\right)^n + 1} \\
  &= \dfrac{1}{\dfrac{1-c}{c} \left[\left(1-\dfrac{a}{n}\right)^{-n/a}\right]^{-a} + 1} \\
  &\xrightarrow{n\to\infty} \dfrac{1}{\dfrac{1-c}{c}e^{-a}+1}
\end{align}
だから整理して$g(a)=\lim_{n\to\infty}f_n(a)$は
\begin{align}
  g(a) &= \dfrac{c}{c + (1-c)e^{-a}} \label{eq:4}
\end{align}

\paragraph{(3)}

$c=2,1,\dfrac{1}{4}$について個別に検討する．いずれも$0<x$に注意する．

\subparagraph{$c=2$の時}

$c=2$を\cref{eq:4}に代入して
\begin{align}
  g(x) = \dfrac{2}{2-e^{-x}}
\end{align}
だから一階・二階微分は
\begin{align}
  g'(x) = \dfrac{-2e^{-x}}{(2-e^{-x})^2} < 0 \\
  g''(x) = \dfrac{2e^{-x}(e^{-x}+2)}{(2-e^{-x})^3} > 0
\end{align}
となって，$g(x)$ は単調減少かつ下に凸である．
また極限値は
\begin{align}
  g(x) \xrightarrow{x \to +0} 2, \quad g(x) \xrightarrow{x \to +\infty} 1
\end{align}
で$g(x)$の増減表は以下のようになる．

\begin{table}[htb]
  \centering
  \caption{$c=2$ のときの $g(x)$ の増減表}
  \label{tab:gx_c2_monotonicity}
  \renewcommand{\arraystretch}{1.4}
  \begin{tabular}{|c|c|c|c|}
    \hline
    $x$   & $0$ & $\cdots$     & $(\infty)$ \\ \hline
    $g'$  &     & $-$          &            \\ \hline
    $g''$ &     & $+$          &            \\ \hline
    $g$   & $(2)$ & $\searrow$（下に凸） & $(1)$ \\ \hline
  \end{tabular}
\end{table}

\subparagraph{$c=1$の時}

$c=1$を\cref{eq:4}に代入して
\begin{align}
  g(x) = 1
\end{align}
である．

\subparagraph{$c=\dfrac{1}{4}$ の時}

$c=\dfrac{1}{4}$を\cref{eq:4}に代入して
\begin{align}
  g(x) = \dfrac{1}{3e^{-x}+1}
\end{align}
だから一階・二階微分は
\begin{align}
  g'(x) = \dfrac{3e^{-x}}{(3e^{-x}+1)^2} > 0 \\
  g''(x) = \dfrac{3e^{-x}(3e^{-x}-1)}{(3e^{-x}+1)^3}
\end{align}
であり，極限値は
\begin{align}
  g(x) \xrightarrow{x \to +0} \dfrac{1}{4}, \quad g(x) \xrightarrow{x \to +\infty} 1
\end{align}
である．従って増減表は以下のようになる．

\begin{table}[htb]
  \centering
  \caption{$C=\dfrac{1}{4}$の時の$g(x)$の増減表}
  \begin{tabular}{|c|c|c|c|c|}
    \hline
    $x$ & $0$ & $\cdots$ & $\log 3$ & $\cdots$ \\
    \hline
    $g'$ & & $+$ & $+$ & $+$ \\
    \hline
    $g''$ & & $+$ & $0$ & $-$ \\
    \hline
    $g$ & & $\nearrow$ & $\dfrac{1}{2}$ & $\nearrow$ \\
    \hline
  \end{tabular}
\end{table}
\casesend

以上からグラフは下図のようになる．

\begin{figure}[htb]
  \centering
  \begin{tikzpicture}[x=1.8cm, y=1.8cm, >=stealth, every node/.style={font=\small}]

    % --- 1. 座標定義 ---
    % ln(3) ≈ 1.0986
    \def\lnThree{1.0986}
    \coordinate (O)       at (0, 0);
    \coordinate (InfPt)   at (\lnThree, 0.5);

    % --- 2. 漸近線 y = 1 (破線) ---
    \draw[dashed, red!80!black, thick] (0, 1.0) -- (4.3, 1.0)
    node[right=2pt, font=\scriptsize, red!90!black] {漸近線 $y = 1$};

    % --- 3. 座標軸 ---
    \draw[->, thick, gray!80!black] (-0.3, 0) -- (4.5, 0) node[right] {$x$};
    \draw[->, thick, gray!80!black] (0, -0.2) -- (0, 2.3) node[above] {$y$};
    \node[below left=2pt] at (O) {$O$};

    % --- 4. 曲線群のプロット (x > 0) ---
    % (i) c = 2: g(x) = 2 / (2 - e^(-x)) (単調減少・下に凸)
    \draw[line width=1.6pt, blue!85!black, domain=0.0:4.2, samples=80]
    plot (\x, {2.0 / (2.0 - exp(-\x))});
    \node[blue!90!black, font=\footnotesize\bfseries, above right] at (1.5, 1.15)
    {$c = 2 \quad \left(g(x) = \frac{2}{2 - e^{-x}}\right)$};

    % (ii) c = 1: g(x) = 1 (定数関数)
    \draw[line width=1.4pt, teal!85!black] (0, 1.0) -- (4.2, 1.0);
    \node[teal!90!black, font=\footnotesize\bfseries, below right] at (2.6, 0.98)
    {$c = 1 \ (g(x) = 1)$};

    % (iii) c = 1/4: g(x) = 1 / (3e^(-x) + 1) (S字型・ロジスティック曲線)
    \draw[line width=1.6pt, purple!85!black, domain=0.0:4.2, samples=80]
    plot (\x, {1.0 / (3.0*exp(-\x) + 1.0)});
    \node[purple!90!black, font=\footnotesize\bfseries, below right] at (1.7, 0.45)
    {$c = \frac{1}{4} \quad \left(g(x) = \frac{1}{3e^{-x} + 1}\right)$};

    % --- 5. 切片・変曲点のプロットと投影破線 ---
    % y 軸切片 (x -> +0)
    % ※ 定義域が x > 0 のため白抜き丸で端点を表示
    \draw[thick, blue!85!black, fill=white] (0, 2.0)  circle (1.4pt);
    \draw[thick, teal!85!black, fill=white] (0, 1.0)  circle (1.4pt);
    \draw[thick, purple!85!black, fill=white] (0, 0.25) circle (1.4pt);

    \node[left=2pt, font=\scriptsize, blue!90!black]   at (0, 2.0)  {$2$};
    \node[left=2pt, font=\scriptsize, red!90!black]    at (0, 1.0)  {$1$};
    \node[left=2pt, font=\scriptsize, purple!90!black] at (0, 0.25) {$\frac{1}{4}$};

    % 変曲点 (log 3, 1/2) (c = 1/4)
    \draw[dashed, gray!60] (\lnThree, 0) -- (InfPt) -- (0, 0.5);
    \fill[purple!90!black] (InfPt) circle (1.5pt);
    \node[above left=1pt, font=\scriptsize, purple!90!black] at (InfPt) {変曲点 $\left(\log 3, \frac{1}{2}\right)$};

    \node[below=2pt, font=\scriptsize] at (\lnThree, 0) {$\log 3$};
    \node[left=2pt, font=\scriptsize]  at (0, 0.5) {$\frac{1}{2}$};

  \end{tikzpicture}
  \caption{$c = 2, 1, \dfrac{1}{4}$ における極限関数 $g(x)$ のグラフ}
  \label{fig:logistic_limit_functions}
\end{figure}

\end{document}
